\documentclass[10pt]{article}
\usepackage{amsmath, amssymb, amsthm}
\usepackage{geometry}
\geometry{a4paper, margin=1in}

\usepackage{lmodern}
\usepackage{parskip}
\usepackage{enumitem}
\usepackage{hyperref}
\usepackage{xcolor}
\usepackage{fancyhdr}
\usepackage{graphicx}

\definecolor{mypink}{RGB}{255, 105, 180}
\definecolor{mydarkpink}{RGB}{199, 21, 133}
\definecolor{mygray}{RGB}{80, 80, 80}
\hypersetup{
    colorlinks=true,
    linkcolor=mydarkpink,
    urlcolor=mydarkpink,
    citecolor=mydarkpink
}

\pagestyle{fancy}
\fancyhf{}
\fancyhead[L]{\textcolor{mygray}{Elijah's Math Notes}}
\fancyhead[R]{\textcolor{mygray}{\thepage}}
\fancyfoot[C]{\textcolor{mygray}{Prepared by Elijah}}

\usepackage{titling}
\pretitle{\begin{center}\Large\color{mypink}}
\posttitle{\par\end{center}\vskip 0.5em}
\preauthor{\begin{center}\large\color{mygray}}
\postauthor{\par\end{center}}
\predate{\begin{center}\small\color{mygray}}
\postdate{\par\end{center}}

\theoremstyle{plain}
\newtheorem{theorem}{Theorem}[section]
\newtheorem{lemma}[theorem]{Lemma}
\newtheorem{proposition}[theorem]{Proposition}
\newtheorem{corollary}[theorem]{Corollary}

\theoremstyle{definition}
\newtheorem{definition}[theorem]{Definition}
\newtheorem{example}[theorem]{Example}

\theoremstyle{remark}
\newtheorem{remark}[theorem]{Remark}
\newtheorem{note}[theorem]{Note}

\newcommand{\R}{\mathbb{R}}
\newcommand{\N}{\mathbb{N}}
\newcommand{\Z}{\mathbb{Z}}
\newcommand{\todo}[1]{\textcolor{mydarkpink}{\textbf{TODO:} #1}}
\newcommand{\highlight}[1]{\textcolor{mypink}{#1}}

\title{Combinatorics and Probability}
\author{Elijah Renner}
\date{\today}

\begin{document}

\maketitle
\tableofcontents

\section{Binomial Coefficients}
 
In combinatorics, \textit{n} choose \textit{k}, denoted as \( \binom{n}{k} \), is the number of ways to choose \textit{k} elements from a set of \textit{n} distinct elements without regard to the order of selection. This concept is closely related to Pascal's Triangle and has applications in binomial expansions, probability theory, and other areas of mathematics.\\

The formula for \( n \choose k \) is given by:
\[
\binom{n}{k} = \frac{n!}{k!(n-k)!}
\]
where \( n! \) represents the factorial of \( n \) and is defined as the product of all positive integers from 1 to \( n \). This expression counts the number of ways to select \textit{k} elements from a set of \textit{n} elements.\\

For example, \( \binom{5}{2} \) is the number of ways to choose 2 elements from a set of 5, which is:
\[
\binom{5}{2} = \frac{5!}{2!(5-2)!} = \frac{5 \times 4}{2 \times 1} = 10
\]

\section{Pascal's Triangle}
Pascal's Triangle is a triangular array of binomial coefficients. Each entry in the triangle is equal to the sum of the two entries directly above it. The general structure of Pascal's Triangle is as follows:

\[
\begin{array}{cccccccc}
    &  &  & 1 &  &  &  &  \\
    &  & 1 &  & 1 &  &  &  \\
    & 1 &  & 2 &  & 1 &  &  \\
    1 &  & 3 &  & 3 &  & 1 &  \\
    \vdots & & \vdots &  & \vdots & & \vdots  \\
\end{array}
\]
Each entry \( \binom{n}{k} \) in Pascal's Triangle represents the number of ways to choose \textit{k} elements from a set of \textit{n} elements. The top row corresponds to \( n = 0 \), the next row to \( n = 1 \), and so on.

\section{Binomial Expansion}
The binomial theorem provides a formula for expanding expressions of the form \( (a + b)^n \). Using \( n \choose k \), the binomial expansion of \( (a + b)^n \) is given by:
\[
(a + b)^n = \sum_{k=0}^{n} \binom{n}{k} a^{n-k} b^k
\]
This formula expresses the powers of a binomial (i.e., \( a + b \)) as a sum of terms involving binomial coefficients. For example, expanding \( (a + b)^3 \) using the binomial theorem gives:
\[
(a + b)^3 = \binom{3}{0} a^3 b^0 + \binom{3}{1} a^2 b^1 + \binom{3}{2} a^1 b^2 + \binom{3}{3} a^0 b^3
\]
Simplifying the terms:
\[
(a + b)^3 = 1 \cdot a^3 + 3 \cdot a^2b + 3 \cdot ab^2 + 1 \cdot b^3 = a^3 + 3a^2b + 3ab^2 + b^3
\]
This result matches what we would expect from multiplying out the terms directly.

\section{Application of Binomial Coefficients in Probability}

To compute the probability that two randomly chosen elements from a set belong to the same group, where the set is partitioned into two groups containing $B$ and $R$ elements, we use binomial coefficients to count the possible outcomes.

\subsection{Total Outcomes}
The total number of ways to choose two elements from a set of $B + R$ elements is given by:
\[
\binom{B + R}{2} = \frac{(B + R)(B + R - 1)}{2}
\]

\subsection{Favorable Outcomes}
There are two favorable cases to consider:
\begin{itemize}
    \item Choosing two elements from the group of size $B$:
    \[
    \binom{B}{2} = \frac{B(B - 1)}{2}
    \]
    \item Choosing two elements from the group of size $R$:
    \[
    \binom{R}{2} = \frac{R(R - 1)}{2}
    \]
\end{itemize}

\subsection{Probability Calculation}
The probability that the two randomly chosen elements belong to the same group is the sum of the probabilities of choosing two elements from either group:
\[
P(\text{same group}) = \frac{\binom{B}{2} + \binom{R}{2}}{\binom{B + R}{2}}
\]


\section{Permutations vs. Combinations}

In combinatorics, permutations and combinations are fundamental concepts used to count the number of ways to arrange or select items from a set.

\subsection{Permutations}

A \textbf{permutation} is an arrangement of objects where the order is important. The number of permutations of \( n \) distinct objects taken \( k \) at a time is given by:
\[
P(n, k) = n \times (n - 1) \times \dots \times (n - k + 1) = \frac{n!}{(n - k)!}
\]
This formula counts the number of ways to arrange \( k \) objects out of \( n \) distinct objects.

\subsection{Combinations}

A \textbf{combination} is a selection of objects where the order does not matter. The number of combinations of \( n \) distinct objects taken \( k \) at a time is given by the binomial coefficient:
\[
C(n, k) = \binom{n}{k} = \frac{n!}{k!(n - k)!}
\]
This formula counts the number of ways to choose \( k \) objects from \( n \) without regard to the order.

\subsection{Key Differences}

The main difference between permutations and combinations is whether the order of selection matters:
\begin{itemize}
    \item \textbf{Permutations}: Order matters.
    \item \textbf{Combinations}: Order does not matter.
\end{itemize}

\subsection{Relationship Between Permutations and Combinations}

Permutations and combinations are related through the following formula:
\[
P(n, k) = C(n, k) \times k!
\]
This relationship shows that the number of permutations is equal to the number of combinations multiplied by \( k! \), accounting for all the possible orders of the \( k \) selected items.

\end{document}